\documentclass[exercice,a4]{thedocument}%
\usepackage[french]{babel}%
\usepackage{pstricks}%
\startdatakeys
    \source{21GENMATAA2}
    \BO{2026}
    \cycle{4}
    \competencesDuSocle{CA,RA}
    \objectifsApprentissage{C4B2eb,C4B3ca}
    \motsCles{Pythagore,Thalès,réciproque de Pythagore}
\enddatakeys
\begin{document}%

\medskip

\begin{minipage}{0.45\linewidth}
Sur la figure ci-contre :

\begin{itemize}
\item[$\bullet~~$]le triangle DCB est rectangle en C ;
\item[$\bullet~~$]les points A, C et B sont alignés ;
\item[$\bullet~~$]les points D, E et B sont alignés ;
\item[$\bullet~~$]AC = 3,2 cm;
\item[$\bullet~~$]CB = 6,8 cm ;
\item[$\bullet~~$]BD = 8,5 cm;
\item[$\bullet~~$]BE = 5,8 cm.
\end{itemize}
\end{minipage}\hfill
\begin{minipage}{0.5\linewidth}
\psset{unit=1cm}
\begin{pspicture}(7.4,7)
\uput[u](3.7,6){La figure n'est pas à l'échelle}
\pspolygon(0.3,1.4)(7,0.4)(2.7,5.4)%ABD
\psline(2.7,5.4)(2.1,1.1)%DC
\psline(1.6,0.2)(4.6,5.3)%CE
\rput{-13}(2.1,1.17){\psframe(0.25,0.25)}
\uput[l](0.3,1.4){A} \uput[dr](7,0.4){B} \uput[dr](2.2,1){C} \uput[u](2.7,5.4){D} \uput[r](3.9,4){E}
\end{pspicture}
\end{minipage}

\medskip

\begin{enumerate}
\item Démontrer que la longueur DC est égale à $5,1$ cm.
\item Calculer l’aire du triangle DCB en cm$^2$.
\item À l’aide de la calculatrice, calculer une valeur approchée de la mesure de l’angle $\widehat{\text{ADC}}$, au degré près.
\item Les droites (AD) et (CE) sont-elles parallèles ?
\end{enumerate}

\bigskip

\showthesource
\end{document}
